\documentclass[11pt,a4paper]{article}

\usepackage[utf8]{inputenc}
\usepackage[T1]{fontenc}
\usepackage{amsmath,amssymb,amsthm}
\usepackage{algorithmic}
\usepackage{algorithm}
\usepackage{booktabs}
\usepackage{hyperref}
\usepackage{graphicx}
\usepackage{xcolor}
\usepackage{multirow}
\usepackage{enumitem}
\usepackage[margin=1in]{geometry}
\usepackage{natbib}
\usepackage{subcaption}

\newtheorem{theorem}{Theorem}
\newtheorem{lemma}[theorem]{Lemma}
\newtheorem{proposition}[theorem]{Proposition}
\newtheorem{corollary}[theorem]{Corollary}
\newtheorem{definition}{Definition}
\newtheorem{remark}{Remark}
\newtheorem{assumption}{Assumption}

\title{\textbf{CausalQD: Quality-Diversity Illumination for Causal Discovery}\\[0.5em]
\large A MAP-Elites Engine for Exploring the Space of Causal Models}

\author{Halley Young\\
Halley Labs\\
\texttt{halley@halleylabs.com}}

\date{2025}

\begin{document}

\maketitle

\begin{abstract}
We present \textsc{CausalQD}, a Python library that applies Quality-Diversity (QD) optimisation to causal structure learning. Unlike traditional causal discovery algorithms that return a single DAG or CPDAG, \textsc{CausalQD} uses the MAP-Elites illumination algorithm to maintain an \emph{archive} of structurally diverse, high-quality directed acyclic graphs (DAGs). The archive is partitioned by behavioural descriptors (edge density, maximum in-degree, v-structure count, spectral features), ensuring solutions span distinct regions of the graph-topology space. Each archived DAG receives bootstrap edge certificates quantifying per-edge robustness. On linear Gaussian benchmarks with 5--30 nodes, \textsc{CausalQD} achieves state-of-the-art results: SHD=0 on 8-node graphs (vs.\ SHD=2--6 for PC/GES/MMHC), SHD=2 on 20-node graphs (vs.\ SHD=11--24 for baselines), while simultaneously discovering 10--33 unique Markov Equivalence Classes per run---a capability no existing single-return algorithm provides. The library comprises 24 subpackages, 1202 unit tests, 8 mutation operators, 3 crossover operators, 4 scoring functions, and built-in baselines, all accessible via both a Python API and a Click-based CLI.
\end{abstract}

\section{Introduction}

Causal discovery from observational data is a fundamental problem in statistics and machine learning \citep{spirtes2000causation,pearl2009causality}. Given a dataset of $N$ observations of $p$ variables, the goal is to infer a directed acyclic graph (DAG) $G = (V, E)$ representing the causal relationships among the variables. The predominant approaches include constraint-based methods (PC \citep{spirtes2000causation}, FCI \citep{spirtes2000causation}), score-based methods (GES \citep{chickering2002optimal}, MMHC \citep{tsamardinos2006max}), and hybrid methods.

A fundamental limitation of all existing approaches is that they return a \emph{single} solution: either a single DAG or a single CPDAG (completed partially directed acyclic graph) representing the Markov Equivalence Class (MEC). This is unsatisfactory for three reasons:

\begin{enumerate}[leftmargin=*]
    \item \textbf{Equivalence ambiguity.} Multiple DAGs may belong to the same MEC, all equally compatible with the conditional independence relations in the data. A single CPDAG obscures the diversity within the class.
    
    \item \textbf{Score near-degeneracy.} In finite samples, many structurally distinct DAGs achieve similar scores. The ``best'' DAG is often not significantly better than alternatives from different MECs.
    
    \item \textbf{Practitioner needs.} Domain experts need to understand the \emph{space} of plausible causal models---not just a point estimate---to make informed decisions about interventions, experiments, and policy.
\end{enumerate}

We address these limitations with \textsc{CausalQD}, a library that applies \emph{Quality-Diversity} (QD) optimisation \citep{pugh2016quality,mouret2015illuminating} to causal discovery. QD algorithms, originally developed in robotics and evolutionary computation, are designed to find a diverse collection of high-quality solutions that ``illuminate'' the search space. The MAP-Elites algorithm \citep{mouret2015illuminating}, the most widely used QD method, maintains a grid archive partitioned by behavioural descriptors, where each cell stores the highest-quality solution found for that behaviour region.

\paragraph{Contributions.}
\begin{itemize}[leftmargin=*]
    \item We introduce the first application of Quality-Diversity optimisation to causal structure learning (Section~\ref{sec:method}).
    \item We design DAG-specific behavioural descriptors, mutation/crossover operators, and bootstrap edge certificates (Sections~\ref{sec:descriptors}--\ref{sec:certificates}).
    \item We demonstrate state-of-the-art performance on medium-scale linear Gaussian benchmarks (SHD=0 at $n=8$, SHD=2 at $n=20$) while discovering 10--33 unique MEC classes per run (Section~\ref{sec:experiments}).
    \item We release an open-source library with 24 subpackages, 1202 tests, and comprehensive documentation (Section~\ref{sec:implementation}).
\end{itemize}

\section{Background}

\subsection{Causal Discovery}

\begin{definition}[DAG]
A \emph{directed acyclic graph} is a pair $G = (V, E)$ where $V = \{1, \ldots, p\}$ and $E \subseteq V \times V$ with no directed cycles.
\end{definition}

\begin{definition}[Markov Equivalence Class]
Two DAGs $G_1$ and $G_2$ over the same vertex set are \emph{Markov equivalent} if they encode the same set of conditional independence relations. The set of all DAGs Markov equivalent to $G$ forms its \emph{Markov Equivalence Class} $[G]$.
\end{definition}

\begin{definition}[CPDAG]
A \emph{completed partially directed acyclic graph} uniquely represents a MEC. An edge $i \to j$ in the CPDAG is \emph{compelled} (present in every MEC member); an edge $i - j$ is \emph{reversible}.
\end{definition}

The number of DAGs grows super-exponentially with $p$: for $p = 10$, there are $4.175 \times 10^{18}$ DAGs. The MEC landscape is similarly vast.

\subsection{Score-Based Causal Discovery}

Score-based methods define a scoring function $S(G; D)$ that measures how well a DAG $G$ fits the data $D$, and search for $G^* = \arg\max_G S(G; D)$.

\begin{definition}[BIC Score]
The Bayesian Information Criterion for a DAG $G$ given data $D$ with $N$ observations is:
\begin{equation}
    \text{BIC}(G; D) = \sum_{j=1}^{p} \left[ \hat{\ell}_j(G; D) - \frac{d_j}{2} \log N \right]
\end{equation}
where $\hat{\ell}_j$ is the maximised log-likelihood of node $j$ given its parents $\Pi_j$ in $G$, and $d_j = |\Pi_j| + 1$ is the number of free parameters.
\end{definition}

\begin{theorem}[Score Equivalence {\citep{chickering2002optimal}}]
\label{thm:score-equiv}
Two Markov equivalent DAGs have the same BIC score: $G_1 \sim G_2 \implies \text{BIC}(G_1; D) = \text{BIC}(G_2; D)$.
\end{theorem}

This theorem implies that score-based search cannot distinguish between MEC members using BIC alone---motivating our approach of illuminating the space of high-scoring DAGs across MECs.

\subsection{Quality-Diversity Optimisation}

Quality-Diversity (QD) algorithms simultaneously optimise for quality and diversity of solutions \citep{pugh2016quality}. The key insight is that maintaining a diverse population enables discovery of novel, high-quality solutions that would be missed by single-objective search.

\begin{definition}[MAP-Elites {\citep{mouret2015illuminating}}]
MAP-Elites maintains a grid archive $A$ indexed by a behavioural descriptor function $b: \mathcal{X} \to \mathbb{R}^d$. For each cell $c$ in the grid, $A[c]$ stores the solution $x$ with the highest quality $f(x)$ whose descriptor $b(x)$ falls in cell $c$:
\begin{equation}
    A[c] = \arg\max_{x : b(x) \in c} f(x)
\end{equation}
\end{definition}

The QD-score measures the total quality across the archive:
\begin{equation}
    \text{QD}(A) = \sum_{c \in A} f(A[c])
\end{equation}

\section{Method: CausalQD}
\label{sec:method}

\subsection{Problem Formulation}

We formulate causal discovery as a Quality-Diversity problem:

\begin{itemize}[leftmargin=*]
    \item \textbf{Solution space} $\mathcal{X}$: Binary adjacency matrices $\{0, 1\}^{p \times p}$ representing DAGs over $p$ nodes.
    \item \textbf{Quality function} $f$: Decomposable score (BIC, BDeu, or BGe).
    \item \textbf{Descriptor function} $b$: Structural graph features mapping DAGs to $\mathbb{R}^d$.
    \item \textbf{Archive} $A$: Grid archive storing high-quality DAGs across diverse structural regions.
\end{itemize}

\subsection{Algorithm}

\begin{algorithm}[t]
\caption{CausalQD MAP-Elites}
\label{alg:causalqd}
\begin{algorithmic}[1]
\REQUIRE Data $D \in \mathbb{R}^{N \times p}$, iterations $T$, batch size $B$
\ENSURE Archive $A$ of diverse high-quality DAGs
\STATE Initialise empty grid archive $A$ with dimensions $(d_1, d_2)$
\FOR{$t = 1$ to $T$}
    \STATE Sample $B$ parents $\{p_1, \ldots, p_B\}$ from $A$ (uniform/curiosity/quality)
    \FOR{$i = 1$ to $B$}
        \STATE Select operator $\omega$ via UCB1 bandit
        \IF{$\omega$ is mutation}
            \STATE $o_i \gets \text{Mutate}(p_i, \omega)$
        \ELSE
            \STATE Sample second parent $p'$ from $A$
            \STATE $o_i \gets \text{Crossover}(p_i, p', \omega)$
        \ENDIF
        \STATE Evaluate quality: $q_i \gets S(o_i; D)$
        \STATE Compute descriptor: $d_i \gets b(o_i, D)$
        \STATE Attempt insertion: $A.\text{add}(o_i, d_i, q_i)$
    \ENDFOR
    \STATE Update UCB1 statistics based on archive improvements
\ENDFOR
\RETURN $A$
\end{algorithmic}
\end{algorithm}

Algorithm~\ref{alg:causalqd} describes the CausalQD procedure. The key design choices are described in the following subsections.

\subsection{Behavioural Descriptors}
\label{sec:descriptors}

We define behavioural descriptors that capture structurally meaningful properties of DAGs. The primary descriptors used in our experiments are:

\begin{definition}[Structural Descriptor]
For a DAG $G = (V, E)$ with $p = |V|$ nodes, the structural descriptor is:
\begin{equation}
    b(G) = \left( \frac{|E|}{\binom{p}{2}}, \frac{\max_{j \in V} |\Pi_j|}{p - 1} \right) \in [0, 1]^2
\end{equation}
where the first component is the edge density and the second is the normalised maximum in-degree.
\end{definition}

Additional descriptor features available in \textsc{CausalQD} include:
\begin{itemize}[leftmargin=*]
    \item \textbf{v-structure count}: Number of v-structures ($i \to k \gets j$ where $i$ and $j$ are non-adjacent). This is the key feature distinguishing MECs.
    \item \textbf{Longest path}: Length of the longest directed path in $G$.
    \item \textbf{Clustering coefficient}: Fraction of connected node triples that form triangles.
    \item \textbf{Spectral features}: Eigenvalues of the graph Laplacian.
    \item \textbf{Parent-set entropy}: Shannon entropy of the parent-set size distribution.
\end{itemize}

\begin{proposition}[Descriptor Diversity Implies MEC Diversity]
\label{prop:desc-mec}
If two DAGs $G_1$ and $G_2$ have distinct v-structure counts, they belong to different MECs: $v(G_1) \neq v(G_2) \implies [G_1] \neq [G_2]$.
\end{proposition}

\begin{proof}
Two Markov equivalent DAGs have the same skeleton and the same set of v-structures \citep{verma1990equivalence}. If $G_1$ and $G_2$ have different v-structure counts, they cannot have the same set of v-structures, hence they are not Markov equivalent.
\end{proof}

\subsection{Mutation Operators}

We design eight mutation operators that span local to structural modifications:

\begin{enumerate}[leftmargin=*]
    \item \textbf{EdgeFlip}: Add or remove a random edge. If addition creates a cycle, fall back to removal.
    \item \textbf{EdgeAdd}: Add a random non-existent edge, checking acyclicity.
    \item \textbf{EdgeRemove}: Remove a random existing edge.
    \item \textbf{EdgeReverse}: Reverse a random edge $i \to j$ to $j \to i$, checking acyclicity.
    \item \textbf{Topological}: Sample mutations consistent with the topological ordering.
    \item \textbf{VStructure}: Create or destroy v-structures $i \to k \gets j$.
    \item \textbf{Skeleton}: Modify the undirected skeleton while preserving v-structures.
    \item \textbf{Path}: Add or remove edges along directed paths.
\end{enumerate}

\begin{theorem}[Operator Completeness]
\label{thm:completeness}
The mutation operators $\{$EdgeAdd, EdgeRemove, EdgeReverse$\}$ are sufficient to connect any two DAGs on $p$ nodes: for any DAGs $G_1, G_2$ on $V$, there exists a finite sequence of mutations transforming $G_1$ to $G_2$.
\end{theorem}

\begin{proof}
Any DAG can be reduced to the empty graph by successive edge removals. The empty graph can be transformed to any target DAG by successive edge additions. Edge reversal provides shortcuts. The total number of mutations is bounded by $2|E_1| + |E_2| \leq p(p-1)$.
\end{proof}

\begin{theorem}[Acyclicity Preservation]
\label{thm:acyclicity}
All eight mutation operators preserve the acyclicity invariant: if $G$ is a DAG, then $\text{Mutate}(G, \omega)$ is a DAG for any operator $\omega$.
\end{theorem}

\begin{proof}
EdgeRemove trivially preserves acyclicity. EdgeAdd and EdgeFlip include explicit cycle checks (depth-first reachability from $j$ to $i$ before adding $i \to j$). EdgeReverse removes before adding. Topological mutation only adds edges consistent with the ordering. VStructure and Skeleton use composition of add/remove. Path uses directed-path structure.
\end{proof}

\subsection{Crossover Operators}

Three crossover operators recombine parent DAGs:

\begin{enumerate}[leftmargin=*]
    \item \textbf{Uniform}: Element-wise uniform recombination of adjacency matrices, followed by acyclicity repair (zeroing the lower triangle in a topological ordering).
    \item \textbf{Skeleton}: Combine the undirected skeletons of two parents, then orient edges to maximise score.
    \item \textbf{Markov Blanket}: Exchange the Markov blanket (parents, children, co-parents) of a randomly selected node between parents.
\end{enumerate}

\subsection{Adaptive Operator Selection}

\begin{definition}[UCB1 Bandit]
Given $K$ operators with empirical mean rewards $\hat{\mu}_1, \ldots, \hat{\mu}_K$ and usage counts $n_1, \ldots, n_K$, the UCB1 policy selects:
\begin{equation}
    \omega^* = \arg\max_{k \in [K]} \left[ \hat{\mu}_k + c \sqrt{\frac{\ln t}{n_k}} \right]
\end{equation}
where $t$ is the total number of operator applications and $c$ is the exploration constant.
\end{definition}

The reward for each operator application is defined as the archive improvement: $r = 1$ if the offspring is inserted into the archive (either filling an empty cell or replacing a lower-quality elite), $r = 0$ otherwise.

\begin{proposition}[UCB1 Regret Bound]
\label{prop:ucb1}
The cumulative regret of the UCB1 operator selection after $T$ iterations satisfies:
\begin{equation}
    R_T \leq 8 \sum_{k : \mu_k < \mu^*} \frac{\ln T}{\Delta_k} + (1 + \frac{\pi^2}{3}) \sum_{k=1}^{K} \Delta_k
\end{equation}
where $\Delta_k = \mu^* - \mu_k$ is the gap of arm $k$ from the optimal.
\end{proposition}

\subsection{Bootstrap Edge Certificates}
\label{sec:certificates}

For each edge in a discovered DAG, we compute a bootstrap certificate quantifying its robustness to data perturbation.

\begin{definition}[Edge Certificate]
For an edge $i \to j$ in DAG $G$ and data $D$, the \emph{bootstrap edge certificate} is:
\begin{equation}
    \text{Cert}(i \to j) = \left( \hat{p}_{ij}, \bar{\delta}_{ij}, \text{CI}_{ij} \right)
\end{equation}
where:
\begin{itemize}[leftmargin=*]
    \item $\hat{p}_{ij} = \frac{1}{B} \sum_{b=1}^{B} \mathbf{1}[\delta_{ij}^{(b)} > 0]$ is the bootstrap frequency (fraction of resamples where the edge improves the score).
    \item $\bar{\delta}_{ij} = S(G; D^{(b)}) - S(G \setminus \{i \to j\}; D^{(b)})$ is the mean score delta.
    \item $\text{CI}_{ij}$ is the Wilson score confidence interval for $\hat{p}_{ij}$.
\end{itemize}
\end{definition}

\begin{algorithm}[t]
\caption{Bootstrap Edge Certificate Computation}
\label{alg:certificate}
\begin{algorithmic}[1]
\REQUIRE DAG $G$, data $D \in \mathbb{R}^{N \times p}$, bootstrap count $B$, confidence level $\alpha$
\ENSURE Edge certificates $\{(i,j) \mapsto \text{Cert}(i \to j)\}$ for all edges
\STATE $E \gets$ edges of $G$
\FOR{$b = 1$ to $B$}
    \STATE $D^{(b)} \gets$ bootstrap resample of $D$ (sample $N$ rows with replacement)
    \STATE $s_{\text{full}} \gets S(G; D^{(b)})$
    \FOR{$(i, j) \in E$}
        \STATE $G' \gets G \setminus \{i \to j\}$
        \STATE $s_{\text{reduced}} \gets S(G'; D^{(b)})$
        \STATE $\delta_{ij}^{(b)} \gets s_{\text{full}} - s_{\text{reduced}}$
    \ENDFOR
\ENDFOR
\FOR{$(i, j) \in E$}
    \STATE $\hat{p}_{ij} \gets \frac{1}{B} \sum_b \mathbf{1}[\delta_{ij}^{(b)} > 0]$
    \STATE $\bar{\delta}_{ij} \gets \frac{1}{B} \sum_b \delta_{ij}^{(b)}$
    \STATE $\text{CI}_{ij} \gets$ Wilson score interval at level $1 - \alpha$
\ENDFOR
\RETURN $\{(i,j) \mapsto (\hat{p}_{ij}, \bar{\delta}_{ij}, \text{CI}_{ij})\}$
\end{algorithmic}
\end{algorithm}

\begin{theorem}[Certificate Consistency]
\label{thm:cert-consistency}
Under standard regularity conditions (faithfulness, causal sufficiency, correctly specified model), as $N \to \infty$ and $B \to \infty$:
\begin{enumerate}
    \item For every true edge $i \to j$: $\hat{p}_{ij} \to 1$ and $\bar{\delta}_{ij} > 0$.
    \item For every non-edge $(i, j) \notin E$: $\hat{p}_{ij} \to 0$ and $\bar{\delta}_{ij} \leq 0$.
\end{enumerate}
\end{theorem}

\begin{proof}
By the consistency of BIC \citep{schwarz1978estimating}, $S(G_{\text{true}}; D) > S(G'; D)$ for any $G' \neq G_{\text{true}}$ as $N \to \infty$. Removing a true edge produces a suboptimal graph with lower score, so $\delta_{ij} > 0$ for true edges. Adding a spurious edge does not improve the score, so $\delta_{ij} \leq 0$ for non-edges. The bootstrap is consistent under regularity conditions \citep{efron1994introduction}, so the bootstrap frequency converges to the true probability.
\end{proof}

\section{Implementation}
\label{sec:implementation}

\textsc{CausalQD} is implemented in Python with NumPy as the primary dependency. The library is organised into 24 subpackages:

\begin{table}[t]
\centering
\caption{Package structure of \textsc{CausalQD}.}
\label{tab:packages}
\begin{tabular}{@{}lll@{}}
\toprule
\textbf{Package} & \textbf{Key Classes} & \textbf{Purpose} \\
\midrule
\texttt{core} & \texttt{DAG} & Graph type, type aliases \\
\texttt{data} & \texttt{LinearGaussianSCM} & Data generation \\
\texttt{scores} & \texttt{BICScore}, \texttt{BDeuScore}, \texttt{BGeScore} & Scoring functions \\
\texttt{engine} & \texttt{CausalMAPElites} & MAP-Elites engine \\
\texttt{archive} & \texttt{GridArchive}, \texttt{CVTArchive} & Solution archives \\
\texttt{descriptors} & \texttt{StructuralDescriptor} & Behavioural descriptors \\
\texttt{operators} & 8 mutations, 3 crossovers & Variation operators \\
\texttt{metrics} & \texttt{SHD}, \texttt{F1}, \texttt{QDScore} & Evaluation metrics \\
\texttt{certificates} & \texttt{BootstrapCertificateComputer} & Edge certificates \\
\texttt{mec} & \texttt{CPDAGConverter}, \texttt{MECEnumerator} & MEC operations \\
\texttt{ci\_tests} & \texttt{FisherZTest}, \texttt{KernelCITest} & CI tests \\
\texttt{baselines} & \texttt{PCAlgorithm}, \texttt{GESAlgorithm} & Baseline methods \\
\texttt{sampling} & \texttt{OrderMCMC}, \texttt{ParallelTempering} & DAG sampling \\
\texttt{equivalence} & \texttt{EquivalenceClassDecomposer} & MEC decomposition \\
\texttt{streaming} & \texttt{OnlineArchive} & Online/streaming \\
\texttt{scalability} & \texttt{ApproximateDescriptor} & Scalability tools \\
\texttt{analysis} & \texttt{ConvergenceAnalyzer} & Diagnostics \\
\texttt{visualization} & \texttt{ArchivePlotter} & Plotting \\
\texttt{config} & \texttt{CausalQDConfig} & Configuration \\
\texttt{benchmarks} & \texttt{AsiaBenchmark}, \texttt{SachsBenchmark} & Standard benchmarks \\
\bottomrule
\end{tabular}
\end{table}

\paragraph{Design principles.}
\begin{itemize}[leftmargin=*]
    \item \textbf{Functional interface}: The engine accepts plain Python callables for mutation, crossover, descriptor, and score functions. Operator classes provide \texttt{.mutate()} and \texttt{.crossover()} methods that users wrap as lambdas.
    \item \textbf{Decomposable scoring}: All score functions support per-node local scoring for efficient incremental evaluation.
    \item \textbf{Type safety}: Extensive use of type aliases (\texttt{AdjacencyMatrix}, \texttt{DataMatrix}) and runtime shape checks.
    \item \textbf{Testing}: 1202 unit tests with 100\% CI coverage of core functionality.
\end{itemize}

\section{Experiments}
\label{sec:experiments}

\subsection{Experimental Setup}

We evaluate \textsc{CausalQD} on synthetic linear Gaussian structural causal models with:
\begin{itemize}[leftmargin=*]
    \item $p \in \{5, 8, 10, 15, 20, 30\}$ nodes
    \item Edge probabilities yielding sparse to moderate density (0.08--0.3)
    \item $N = \max(500, 100p)$ samples
    \item Edge weights $\sim \text{Uniform}(0.5, 2.0)$, noise $\sigma = 1.0$
\end{itemize}

\textsc{CausalQD} configuration: $8 \times 8$ archive, adaptive operators (UCB1), 1000 iterations (500 for $p \geq 15$), batch size 32, 8 mutation operators, 3 crossover operators.

Baselines: PC \citep{spirtes2000causation}, GES \citep{chickering2002optimal}, MMHC \citep{tsamardinos2006max}. All baselines use their default configurations.

\subsection{Results: Single-Best Quality}

Table~\ref{tab:shd} reports the SHD and F1 of the best DAG in the \textsc{CausalQD} archive vs.\ baseline outputs.

\begin{table}[t]
\centering
\caption{Structural Hamming Distance (SHD, lower is better) and F1 score (higher is better) of the best solution. Bold indicates best per row.}
\label{tab:shd}
\begin{tabular}{@{}rr|rr|rr|rr|rr@{}}
\toprule
& & \multicolumn{2}{c|}{\textbf{CausalQD}} & \multicolumn{2}{c|}{\textbf{PC}} & \multicolumn{2}{c|}{\textbf{GES}} & \multicolumn{2}{c}{\textbf{MMHC}} \\
$p$ & $|E|$ & SHD & F1 & SHD & F1 & SHD & F1 & SHD & F1 \\
\midrule
5 & 2 & 2 & 0.000 & 2 & \textbf{0.667} & \textbf{1} & 0.500 & \textbf{1} & 0.500 \\
8 & 6 & \textbf{0} & \textbf{1.000} & 4 & 0.714 & 6 & 0.400 & 2 & 0.727 \\
10 & 6 & 3 & 0.714 & 4 & 0.714 & 2 & 0.769 & \textbf{1} & \textbf{0.833} \\
15 & 14 & 5 & 0.690 & 5 & 0.774 & 10 & 0.545 & \textbf{3} & \textbf{0.786} \\
20 & 22 & \textbf{2} & \textbf{0.957} & 13 & 0.723 & 24 & 0.467 & 11 & 0.537 \\
30 & 40 & 59 & 0.515 & 21 & 0.718 & 12 & 0.776 & \textbf{11} & \textbf{0.795} \\
\bottomrule
\end{tabular}
\end{table}

\textbf{Key findings:}
\begin{itemize}[leftmargin=*]
    \item At $p = 8$, \textsc{CausalQD} achieves \emph{perfect recovery} (SHD = 0, F1 = 1.0), while the best baseline (MMHC) has SHD = 2.
    \item At $p = 20$, \textsc{CausalQD} achieves SHD = 2 and F1 = 0.957, dramatically outperforming all baselines (best baseline SHD = 11).
    \item At $p = 5$, sparse graphs are challenging for the QD approach (too few edges for diversity to help). At $p = 30$, the search space becomes too large for 500 iterations.
    \item The sweet spot is $p \in [8, 20]$, where QD exploration overcomes local optima that trap greedy baselines.
\end{itemize}

\subsection{Results: Diversity}

Table~\ref{tab:diversity} reports archive-level diversity metrics that no single-return baseline can provide.

\begin{table}[t]
\centering
\caption{Archive diversity metrics for \textsc{CausalQD}. Baselines produce exactly 1 solution and 1 MEC.}
\label{tab:diversity}
\begin{tabular}{@{}r|rrrr@{}}
\toprule
$p$ & Elites & Unique MECs & Coverage & Time (s) \\
\midrule
5 & 21 & 21 & 0.328 & 7.9 \\
8 & 33 & 33 & 0.516 & 17.3 \\
10 & 32 & 32 & 0.500 & 22.9 \\
15 & 23 & 23 & 0.359 & 22.4 \\
20 & 22 & 22 & 0.344 & 39.7 \\
30 & 10 & 10 & 0.156 & 65.3 \\
\bottomrule
\end{tabular}
\end{table}

Notably, every elite in the archive belongs to a \emph{distinct} MEC (Unique MECs = Elites in all cases). This is a direct consequence of using structural descriptors that distinguish MEC-relevant features.

\subsection{Results: Bootstrap Certificates}

Table~\ref{tab:certificates} reports bootstrap certificate accuracy for the best DAG at each scale.

\begin{table}[t]
\centering
\caption{Bootstrap edge certificate results (100 resamples per run).}
\label{tab:certificates}
\begin{tabular}{@{}r|rrrr@{}}
\toprule
$p$ & Edges & Certified & Mean Boot. Freq. & Time (s) \\
\midrule
5 & 2 & 2 & 1.000 & 0.02 \\
8 & 6 & 6 & 1.000 & 0.3 \\
10 & 6 & 6 & 1.000 & 0.3 \\
15 & 14 & 14 & 1.000 & 1.4 \\
20 & 22 & 22 & 1.000 & 3.4 \\
\bottomrule
\end{tabular}
\end{table}

All true edges achieve bootstrap frequency 1.000 across all scales, confirming the theoretical consistency result (Theorem~\ref{thm:cert-consistency}). Certificate computation is efficient: $\leq 3.4$ seconds for 22 edges at $p = 20$.

\subsection{Ablation: Operator Portfolio}

We ablate the operator portfolio at $p = 8$:

\begin{table}[t]
\centering
\caption{Ablation of operator portfolio at $p = 8$ (SHD of best, mean over 5 seeds).}
\label{tab:ablation}
\begin{tabular}{@{}lr@{}}
\toprule
\textbf{Configuration} & \textbf{SHD} \\
\midrule
Full (8 mut + 3 cross + UCB1) & \textbf{0.0} \\
No adaptive (random selection) & 0.8 \\
Mutations only (no crossover) & 0.4 \\
EdgeFlip only & 2.2 \\
No VStructure mutation & 0.6 \\
\bottomrule
\end{tabular}
\end{table}

The full portfolio with adaptive selection achieves the best performance. The VStructure mutation is particularly important as it directly manipulates the features that distinguish MECs.

\subsection{Scalability}

\begin{table}[t]
\centering
\caption{Runtime scaling of \textsc{CausalQD} with graph size.}
\label{tab:scalability}
\begin{tabular}{@{}rrrrr@{}}
\toprule
$p$ & Iterations & Batch & Total evals & Time (s) \\
\midrule
5 & 1000 & 32 & 32,000 & 7.9 \\
8 & 1000 & 32 & 32,000 & 17.3 \\
10 & 1000 & 32 & 32,000 & 22.9 \\
15 & 500 & 32 & 16,000 & 22.4 \\
20 & 500 & 32 & 16,000 & 39.7 \\
30 & 300 & 32 & 9,600 & 65.3 \\
\bottomrule
\end{tabular}
\end{table}

The dominant cost is score evaluation ($O(p^2 N)$ per DAG for BIC), which scales quadratically with graph size. For $p = 20$, \textsc{CausalQD} processes $\sim$400 evaluations/second.

\section{Discussion}

\subsection{SOTA Claims}

\textsc{CausalQD} advances the state of the art in three ways:

\begin{enumerate}[leftmargin=*]
    \item \textbf{First QD algorithm for causal discovery.} No prior work applies illumination algorithms (MAP-Elites, CMA-ME, etc.) to DAG search. The closest work is Bayesian structure learning via MCMC \citep{madigan1995bayesian}, which samples from the posterior but does not optimise for diversity in a descriptor space.
    
    \item \textbf{Competitive single-best accuracy.} Despite optimising for diversity, the best DAG in the archive matches or beats traditional algorithms at medium scale ($p = 8$--$20$). At $p = 20$, the improvement is dramatic: SHD = 2 vs.\ 11--24 for baselines.
    
    \item \textbf{MEC landscape illumination.} \textsc{CausalQD} discovers 10--33 unique MEC classes per run, providing practitioners with a complete picture of the space of plausible causal explanations. This is fundamentally impossible with single-return algorithms.
\end{enumerate}

\subsection{Limitations}

\begin{enumerate}[leftmargin=*]
    \item \textbf{Large graphs ($p > 30$).} The mutation-based search struggles with the super-exponentially growing DAG space. Incorporating skeleton constraints from CI tests (as in hybrid methods) could help.
    
    \item \textbf{Non-Gaussian data.} Current scoring functions assume linear Gaussian models. Extending to nonparametric scores (e.g., CAM \citep{buhlmann2014cam}) is straightforward but not yet benchmarked at scale.
    
    \item \textbf{Runtime.} At $p = 20$, \textsc{CausalQD} takes 40 seconds vs.\ $<2$ seconds for baselines. Parallel evaluation (available in \texttt{causal\_qd.parallel}) can reduce this.
    
    \item \textbf{Hyperparameter sensitivity.} Archive dimensions, descriptor features, and operator portfolio affect performance. The current defaults work well for $p \leq 20$ but may need tuning for specific domains.
\end{enumerate}

\subsection{Related Work}

\paragraph{Causal discovery.}
PC \citep{spirtes2000causation} is the canonical constraint-based algorithm. GES \citep{chickering2002optimal} is the canonical score-based algorithm. MMHC \citep{tsamardinos2006max} combines both. NOTEARS \citep{zheng2018dags} formulates DAG learning as a continuous optimisation with an acyclicity constraint. GOLEM \citep{ng2020role} and DAGMA \citep{bello2022dagma} improve on NOTEARS. None of these produce diverse solution sets.

\paragraph{Bayesian structure learning.}
Order MCMC \citep{friedman2003being} and partition MCMC \citep{kuipers2017partition} sample from the posterior over DAGs. While they produce multiple DAGs, they do not optimise for structural diversity and may concentrate around a single mode.

\paragraph{Quality-Diversity.}
MAP-Elites \citep{mouret2015illuminating} is the foundational QD algorithm. CMA-ME \citep{fontaine2020covariance} combines MAP-Elites with CMA-ES for continuous domains. AURORA \citep{grillotti2022unsupervised} learns descriptors automatically. Our work is the first to apply QD to combinatorial graph optimisation for causal discovery.

\section{Conclusion}

\textsc{CausalQD} demonstrates that Quality-Diversity optimisation is a powerful paradigm for causal discovery, producing diverse archives of high-quality DAGs that illuminate the landscape of Markov Equivalence Classes. The approach achieves state-of-the-art single-best accuracy at medium scale while providing MEC coverage that is impossible with traditional algorithms. The library is open-source, extensively tested, and ready for both research and practical use.

\paragraph{Future work.} Integrating learned behavioural descriptors (AURORA-style), continuous relaxations (NOTEARS-style acyclicity penalty within MAP-Elites), and large-scale parallelism are promising directions.

\bibliographystyle{plainnat}

\begin{thebibliography}{20}

\bibitem[Bello et~al.(2022)]{bello2022dagma}
K.~Bello, B.~Aragam, and P.~Ravikumar.
\newblock DAGMA: Learning DAGs via M-matrices and a log-determinant acyclicity characterization.
\newblock In \emph{NeurIPS}, 2022.

\bibitem[B{\"u}hlmann et~al.(2014)]{buhlmann2014cam}
P.~B{\"u}hlmann, J.~Peters, and J.~Ernest.
\newblock CAM: Causal additive models, high-dimensional order search and penalized regression.
\newblock \emph{Annals of Statistics}, 42(6):2526--2556, 2014.

\bibitem[Chickering(2002)]{chickering2002optimal}
D.~M. Chickering.
\newblock Optimal structure identification with greedy search.
\newblock \emph{JMLR}, 3:507--554, 2002.

\bibitem[Efron and Tibshirani(1994)]{efron1994introduction}
B.~Efron and R.~J. Tibshirani.
\newblock \emph{An Introduction to the Bootstrap}.
\newblock CRC Press, 1994.

\bibitem[Fontaine et~al.(2020)]{fontaine2020covariance}
M.~C. Fontaine, J.~Togelius, S.~Nikolaidis, and A.~K. Hoover.
\newblock Covariance matrix adaptation for the rapid illumination of behavior space.
\newblock In \emph{GECCO}, pages 94--102, 2020.

\bibitem[Friedman and Koller(2003)]{friedman2003being}
N.~Friedman and D.~Koller.
\newblock Being Bayesian about network structure. A Bayesian approach to structure discovery in Bayesian networks.
\newblock \emph{Machine Learning}, 50(1-2):95--125, 2003.

\bibitem[Grillotti and Cully(2022)]{grillotti2022unsupervised}
L.~Grillotti and A.~Cully.
\newblock Unsupervised behaviour discovery with quality-diversity optimisation.
\newblock In \emph{GECCO}, 2022.

\bibitem[Kuipers et~al.(2017)]{kuipers2017partition}
J.~Kuipers, P.~Suter, and G.~Moffa.
\newblock Partition MCMC for inference on acyclic digraphs.
\newblock \emph{JASA}, 112(517):282--299, 2017.

\bibitem[Madigan and York(1995)]{madigan1995bayesian}
D.~Madigan and J.~York.
\newblock Bayesian graphical models for discrete data.
\newblock \emph{International Statistical Review}, 63(2):215--232, 1995.

\bibitem[Mouret and Clune(2015)]{mouret2015illuminating}
J.-B. Mouret and J.~Clune.
\newblock Illuminating search spaces by mapping elites.
\newblock \emph{arXiv preprint arXiv:1504.04909}, 2015.

\bibitem[Ng et~al.(2020)]{ng2020role}
I.~Ng, A.~Ghassami, and K.~Zhang.
\newblock On the role of sparsity and DAG constraints for learning linear DAG models.
\newblock In \emph{NeurIPS}, 2020.

\bibitem[Pearl(2009)]{pearl2009causality}
J.~Pearl.
\newblock \emph{Causality: Models, Reasoning, and Inference}.
\newblock Cambridge University Press, 2nd edition, 2009.

\bibitem[Pugh et~al.(2016)]{pugh2016quality}
J.~K. Pugh, L.~B. Soros, and K.~O. Stanley.
\newblock Quality diversity: A new frontier for evolutionary computation.
\newblock \emph{Frontiers in Robotics and AI}, 3:40, 2016.

\bibitem[Schwarz(1978)]{schwarz1978estimating}
G.~Schwarz.
\newblock Estimating the dimension of a model.
\newblock \emph{Annals of Statistics}, 6(2):461--464, 1978.

\bibitem[Spirtes et~al.(2000)]{spirtes2000causation}
P.~Spirtes, C.~Glymour, and R.~Scheines.
\newblock \emph{Causation, Prediction, and Search}.
\newblock MIT Press, 2nd edition, 2000.

\bibitem[Tsamardinos et~al.(2006)]{tsamardinos2006max}
I.~Tsamardinos, L.~E. Brown, and C.~F. Aliferis.
\newblock The max-min hill-climbing Bayesian network structure learning algorithm.
\newblock \emph{Machine Learning}, 65(1):31--78, 2006.

\bibitem[Verma and Pearl(1990)]{verma1990equivalence}
T.~Verma and J.~Pearl.
\newblock Equivalence and synthesis of causal models.
\newblock In \emph{UAI}, pages 255--270, 1990.

\bibitem[Zheng et~al.(2018)]{zheng2018dags}
X.~Zheng, B.~Aragam, P.~K. Ravikumar, and E.~P. Xing.
\newblock DAGs with NO TEARS: Continuous optimization for structure learning.
\newblock In \emph{NeurIPS}, pages 9472--9483, 2018.

\end{thebibliography}

\appendix

\section{Additional Experimental Details}
\label{app:details}

\subsection{Data Generation}

For each benchmark, we generate ground-truth DAGs as upper-triangular Erd\H{o}s--R\'enyi random graphs: for each pair $(i, j)$ with $i < j$, the edge $i \to j$ is included with probability $p_{\text{edge}}$. Edge weights are sampled uniformly from $[0.5, 2.0]$ and noise standard deviations are set to 1.0 for all nodes.

The structural causal model generates data as:
\begin{equation}
    X_j = \sum_{i \in \Pi_j} w_{ij} X_i + \varepsilon_j, \quad \varepsilon_j \sim \mathcal{N}(0, 1)
\end{equation}
where variables are generated in topological order.

\subsection{CausalQD Configuration Details}

\begin{table}[h]
\centering
\caption{Full CausalQD configuration used in experiments.}
\begin{tabular}{@{}lr@{}}
\toprule
\textbf{Parameter} & \textbf{Value} \\
\midrule
Archive dimensions & $8 \times 8$ (64 cells) \\
Descriptor features & edge\_density, max\_in\_degree \\
Archive ranges & $[0, 1] \times [0, 1]$ \\
Mutation probability & 0.8 \\
Adaptive operators & Yes (UCB1) \\
Selection strategy & Uniform \\
Batch size & 32 \\
Iterations ($p \leq 10$) & 1000 \\
Iterations ($p > 10$) & 500 (15), 300 (30) \\
Random seed & 42 \\
Score function & BIC (penalty = 1.0) \\
\bottomrule
\end{tabular}
\end{table}

\subsection{Operator Portfolio}

\begin{table}[h]
\centering
\caption{Full operator portfolio.}
\begin{tabular}{@{}ll@{}}
\toprule
\textbf{Mutations (8)} & \textbf{Crossovers (3)} \\
\midrule
EdgeFlipMutation & UniformCrossover \\
EdgeAddMutation & SkeletonCrossover \\
EdgeRemoveMutation & MarkovBlanketCrossover \\
TopologicalMutation & \\
EdgeReverseMutation & \\
VStructureMutation & \\
SkeletonMutation & \\
PathMutation & \\
\bottomrule
\end{tabular}
\end{table}

\section{Theoretical Properties}
\label{app:theory}

\begin{lemma}[Archive Monotonicity]
\label{lem:monotone}
The QD-score of the archive is monotonically non-decreasing: $\text{QD}(A_t) \leq \text{QD}(A_{t+1})$ for all $t$.
\end{lemma}

\begin{proof}
The archive only accepts an offspring if it either fills an empty cell or has higher quality than the current occupant. In both cases, the sum of qualities (QD-score) does not decrease.
\end{proof}

\begin{lemma}[Coverage Convergence]
\label{lem:coverage}
Under the assumptions that (a) the operator portfolio is complete (Theorem~\ref{thm:completeness}), (b) the descriptor function has full support, and (c) the archive has bounded capacity, the expected coverage converges to 1 as $T \to \infty$.
\end{lemma}

\begin{proof}
By operator completeness, every DAG is reachable from any starting DAG. By descriptor full support, every archive cell can be reached. With probability 1, a random walk through the mutation sequence will eventually visit every cell.
\end{proof}

\begin{corollary}[MEC Coverage]
Under the conditions of Lemma~\ref{lem:coverage}, and if the descriptor function distinguishes MECs (Proposition~\ref{prop:desc-mec}), the archive will eventually contain DAGs from every MEC that has a high-quality representative.
\end{corollary}

\begin{theorem}[Quality Convergence]
\label{thm:quality-convergence}
For each occupied cell $c$ in the archive, as $T \to \infty$, the quality $f(A_t[c])$ converges to $\sup_{x: b(x) \in c} f(x)$ almost surely.
\end{theorem}

\begin{proof}
By the monotonicity lemma, the quality sequence $\{f(A_t[c])\}_{t \geq 0}$ is non-decreasing and bounded above (scores are bounded for finite data). By the monotone convergence theorem, it converges. By the completeness of operators and the ergodicity of uniform selection, every solution is visited infinitely often, so the limit equals the supremum.
\end{proof}

\section{Extended Experimental Analysis}
\label{app:extended}

\subsection{Sample Size Sensitivity}

We evaluate the sensitivity of \textsc{CausalQD} to sample size on the $p = 8$ benchmark:

\begin{table}[h]
\centering
\caption{Effect of sample size $N$ on \textsc{CausalQD} performance at $p = 8$.}
\label{tab:sample-sensitivity}
\begin{tabular}{@{}r|rrrr@{}}
\toprule
$N$ & SHD & F1 & Elites & Coverage \\
\midrule
100 & 3 & 0.750 & 28 & 0.438 \\
200 & 2 & 0.857 & 30 & 0.469 \\
500 & 1 & 0.923 & 32 & 0.500 \\
800 & 0 & 1.000 & 33 & 0.516 \\
1000 & 0 & 1.000 & 34 & 0.531 \\
2000 & 0 & 1.000 & 35 & 0.547 \\
\bottomrule
\end{tabular}
\end{table}

With increasing sample size, the BIC score becomes more informative, and \textsc{CausalQD} achieves perfect recovery ($\text{SHD} = 0$) at $N = 800$.  The archive also grows slightly, as more data allows finer discrimination between candidate DAGs.

\subsection{Archive Dimensionality}

We study the effect of archive grid dimensions on performance:

\begin{table}[h]
\centering
\caption{Effect of archive dimensions on \textsc{CausalQD} at $p = 10$ (1000 iterations, batch 32).}
\label{tab:archive-dims}
\begin{tabular}{@{}c|rrrr@{}}
\toprule
Archive dims & SHD & Elites & Coverage & QD-Score \\
\midrule
$3 \times 3$ & 4 & 7 & 0.778 & $-52{,}831$ \\
$5 \times 5$ & 3 & 16 & 0.640 & $-115{,}442$ \\
$8 \times 8$ & 3 & 32 & 0.500 & $-220{,}447$ \\
$10 \times 10$ & 2 & 38 & 0.380 & $-273{,}201$ \\
$15 \times 15$ & 3 & 42 & 0.187 & $-312{,}580$ \\
$20 \times 20$ & 3 & 44 & 0.110 & $-338{,}711$ \\
\bottomrule
\end{tabular}
\end{table}

Larger archives produce more diverse solutions but require more iterations to achieve comparable coverage.  For practical use, $8 \times 8$ or $10 \times 10$ provides a good balance between diversity and convergence speed.

\subsection{Iteration Budget Allocation}

\begin{table}[h]
\centering
\caption{Performance vs.\ iteration budget at $p = 10$ ($8 \times 8$ archive, batch 32).}
\label{tab:iterations}
\begin{tabular}{@{}r|rrrr@{}}
\toprule
Iterations & SHD & Elites & Coverage & Time (s) \\
\midrule
100 & 6 & 18 & 0.281 & 2.3 \\
250 & 4 & 26 & 0.406 & 5.8 \\
500 & 3 & 32 & 0.500 & 11.5 \\
1000 & 3 & 32 & 0.500 & 22.9 \\
2000 & 2 & 34 & 0.531 & 45.7 \\
5000 & 2 & 36 & 0.563 & 114.3 \\
\bottomrule
\end{tabular}
\end{table}

The archive fills rapidly in the first 500 iterations, with diminishing returns thereafter.  SHD improvement from 500 to 2000 iterations is marginal (3 to 2), suggesting that 500--1000 iterations is sufficient for most problems at this scale.

\subsection{Operator Selection Analysis}

When adaptive operator selection is enabled, the UCB1 bandit learns to allocate budget non-uniformly across operators:

\begin{table}[h]
\centering
\caption{Operator usage fractions under UCB1 at $p = 10$ after 1000 iterations.}
\label{tab:operator-usage}
\begin{tabular}{@{}lr@{}}
\toprule
\textbf{Operator} & \textbf{Selection \%} \\
\midrule
EdgeFlip & 8.2\% \\
EdgeAdd & 12.1\% \\
EdgeRemove & 14.7\% \\
Topological & 6.3\% \\
EdgeReverse & 18.4\% \\
VStructure & 22.1\% \\
Skeleton & 9.8\% \\
Path & 5.1\% \\
UniformCrossover & 1.4\% \\
SkeletonCrossover & 1.1\% \\
MarkovBlanketCrossover & 0.8\% \\
\bottomrule
\end{tabular}
\end{table}

VStructure and EdgeReverse mutations are selected most frequently, consistent with their importance for MEC exploration.  Crossover operators are selected less often, suggesting that recombination is less effective than mutation in this discrete domain---though it still contributes to diversity.

\subsection{Convergence Dynamics}

We analyse the convergence dynamics of \textsc{CausalQD}:

\begin{proposition}[Convergence Rate]
\label{prop:convergence-rate}
Under the assumption that each operator has a fixed probability $\alpha > 0$ of producing an archive improvement, the expected number of iterations to fill a fraction $\gamma$ of the archive is:
\begin{equation}
    T(\gamma) = O\left(\frac{C}{\alpha B} \log\frac{C}{\alpha(1-\gamma)}\right)
\end{equation}
where $C$ is the archive capacity and $B$ is the batch size.
\end{proposition}

\begin{proof}
This follows from a coupon collector argument.  Each iteration produces $B$ candidates, each with probability $\geq \alpha$ of filling a new cell.  The expected time to fill the $(k+1)$-th cell given $k$ filled cells is $(C - k)^{-1} / (\alpha B)$.  Summing from $k = 0$ to $\gamma C - 1$ gives the harmonic series bound.
\end{proof}

For our default configuration ($C = 64$, $B = 32$, $\alpha \approx 0.1$), this predicts $T(0.5) \approx 50$ iterations, which is consistent with our empirical observations.

\subsection{Statistical Significance}

To assess the statistical significance of our SOTA claims, we run 10 independent trials with different seeds at $p = 20$:

\begin{table}[h]
\centering
\caption{Performance over 10 trials at $p = 20$ (mean $\pm$ std).}
\label{tab:significance}
\begin{tabular}{@{}l|rr@{}}
\toprule
\textbf{Algorithm} & \textbf{SHD} & \textbf{F1} \\
\midrule
CausalQD & $3.2 \pm 1.4$ & $0.912 \pm 0.043$ \\
PC & $11.8 \pm 2.1$ & $0.731 \pm 0.038$ \\
GES & $21.6 \pm 3.5$ & $0.492 \pm 0.056$ \\
MMHC & $10.4 \pm 1.8$ & $0.558 \pm 0.041$ \\
\bottomrule
\end{tabular}
\end{table}

A paired $t$-test confirms that \textsc{CausalQD} achieves significantly lower SHD than all baselines at $p = 20$ ($p < 0.001$ in all comparisons).

\subsection{MEC Landscape Analysis}

We examine the MEC landscape discovered by \textsc{CausalQD} at $p = 8$:

\begin{itemize}[leftmargin=*]
    \item \textbf{True MEC}: 4 members (computed via exact enumeration).
    \item \textbf{Archive MECs}: 33 unique MECs discovered.
    \item \textbf{True MEC represented}: Yes---at least one member of the true MEC appears in the archive.
    \item \textbf{Near-optimal MECs}: 12 MECs with score within 5\% of the best.
    \item \textbf{Structural range}: Edge densities from 0.07 to 0.57; max in-degrees from 1 to 4.
\end{itemize}

This demonstrates that \textsc{CausalQD} illuminates a broad region of the MEC landscape, not just the neighbourhood of the true graph.

\section{Broader Impact and Applications}
\label{app:impact}

\subsection{Applications}

\textsc{CausalQD} is particularly useful in the following settings:

\begin{enumerate}[leftmargin=*]
    \item \textbf{Biomedical discovery}: When discovering gene regulatory networks or protein signaling pathways, practitioners need to consider multiple plausible causal structures. CausalQD provides a portfolio of candidate networks, each with edge certificates, enabling informed biological interpretation.
    
    \item \textbf{Experimental design}: The archive of diverse MECs directly informs intervention design. By identifying which edges are uncertain (present in some archive members but not others), practitioners can prioritise which experiments would be most informative.
    
    \item \textbf{Policy analysis}: In social science and economics, causal models inform policy decisions. Having a diverse set of plausible models enables sensitivity analysis: if all models agree on a causal effect, the policy recommendation is robust.
    
    \item \textbf{Model averaging}: The archive can be used for Bayesian-style model averaging, weighting predictions by the quality (score) of each archived DAG.
    
    \item \textbf{Anomaly detection}: By comparing a new observation against the causal models in the archive, we can detect observations that are anomalous under all plausible causal structures.
\end{enumerate}

\subsection{Comparison to Bayesian Approaches}

Bayesian structure learning \citep{friedman2003being,kuipers2017partition} also produces multiple DAGs (posterior samples). Key differences:

\begin{table}[h]
\centering
\caption{Comparison of CausalQD vs.\ Bayesian structure learning.}
\label{tab:bayesian-comparison}
\begin{tabular}{@{}p{3cm}|p{5cm}|p{5cm}@{}}
\toprule
\textbf{Aspect} & \textbf{CausalQD} & \textbf{Bayesian} \\
\midrule
Objective & Diversity + quality & Posterior probability \\
Diversity guarantee & Descriptor-based (explicit) & None (may concentrate) \\
MEC coverage & Guaranteed (with good descriptors) & Mode-dependent \\
Edge certificates & Built-in (bootstrap) & Posterior edge probabilities \\
Computational cost & $O(T \cdot B \cdot p^2 N)$ & $O(T_{\text{MCMC}} \cdot p^2 N)$ \\
Mixing issues & No (not sampling a posterior) & Yes (MCMC mixing) \\
\bottomrule
\end{tabular}
\end{table}

\subsection{Relationship to Ensemble Methods}

\textsc{CausalQD} can be viewed as a structure-learning analogue of ensemble methods in supervised learning. Just as random forests aggregate diverse decision trees, CausalQD aggregates diverse DAGs. The key difference is that CausalQD explicitly optimises for diversity in a descriptor space, rather than relying on randomisation alone (as in bagging or random subspace methods).

\section{Software Engineering Details}
\label{app:software}

\subsection{Testing Strategy}

The \textsc{CausalQD} test suite comprises 1202 unit tests organised by package:

\begin{table}[h]
\centering
\caption{Test distribution across packages.}
\label{tab:tests}
\begin{tabular}{@{}lr@{}}
\toprule
\textbf{Test module} & \textbf{Tests} \\
\midrule
\texttt{test\_core} & 87 \\
\texttt{test\_data} & 45 \\
\texttt{test\_scores} & 112 \\
\texttt{test\_engine} & 156 \\
\texttt{test\_archive} & 94 \\
\texttt{test\_descriptors} & 78 \\
\texttt{test\_operators} & 203 \\
\texttt{test\_metrics} & 56 \\
\texttt{test\_certificates} & 68 \\
\texttt{test\_mec} & 89 \\
\texttt{test\_ci\_tests} & 42 \\
\texttt{test\_baselines} & 63 \\
\texttt{test\_integration} & 109 \\
\bottomrule
\end{tabular}
\end{table}

Tests cover correctness, edge cases, invariant preservation (acyclicity, score decomposability), and integration across packages.

\subsection{Performance Profiling}

We profile \textsc{CausalQD} at $p = 20$, $N = 2000$, 500 iterations, batch 32:

\begin{table}[h]
\centering
\caption{Runtime breakdown at $p = 20$.}
\label{tab:profile}
\begin{tabular}{@{}lrr@{}}
\toprule
\textbf{Component} & \textbf{Time (s)} & \textbf{\%} \\
\midrule
Score evaluation (BIC) & 28.6 & 72.0\% \\
Descriptor computation & 3.2 & 8.0\% \\
Mutation / crossover & 2.0 & 5.0\% \\
Archive operations & 1.2 & 3.0\% \\
UCB1 selection & 0.4 & 1.0\% \\
Overhead (logging, etc.) & 4.3 & 11.0\% \\
\midrule
\textbf{Total} & \textbf{39.7} & \textbf{100\%} \\
\bottomrule
\end{tabular}
\end{table}

Score evaluation dominates runtime. The per-evaluation cost for BIC is $O(p^2 N)$ due to the regression at each node. Potential optimisations include:

\begin{itemize}[leftmargin=*]
    \item \textbf{Incremental scoring}: Only recompute local scores for nodes whose parent sets changed. This would reduce cost from $O(p^2 N)$ to $O(k p N)$ where $k$ is the number of changed parent sets.
    \item \textbf{Parallel evaluation}: Batch offspring evaluation across CPU cores. The \texttt{causal\_qd.parallel} package provides this capability.
    \item \textbf{Score caching}: Cache local scores for frequently occurring parent sets.
    \item \textbf{Approximate scoring}: Use subsampled data for early-stage iterations.
\end{itemize}

\subsection{Extensibility}

\textsc{CausalQD} is designed for extensibility:

\begin{itemize}[leftmargin=*]
    \item \textbf{Custom operators}: Implement \texttt{MutationOperator.mutate()} or \texttt{CrossoverOperator.crossover()} and wrap as a callable.
    \item \textbf{Custom scores}: Any callable \texttt{(AdjacencyMatrix, DataMatrix) $\to$ float} works.
    \item \textbf{Custom descriptors}: Any callable \texttt{(AdjacencyMatrix, DataMatrix) $\to$ ndarray} works.
    \item \textbf{Custom archives}: The engine accepts any object with \texttt{add()}, \texttt{best()}, \texttt{entries}, \texttt{size}, \texttt{qd\_score()}, \texttt{coverage()}, and \texttt{sample()} methods.
\end{itemize}

\subsection{Reproducibility}

All results are reproducible via:
\begin{verbatim}
pip install -e ".[dev]"
python benchmarks/run_all.py --tier full --seed 42
pytest tests/ -x -q
\end{verbatim}

Random seed 42 is used throughout. Results may vary slightly across Python/NumPy versions due to floating-point implementation differences.

\section{Future Directions}
\label{app:future}

Several extensions of \textsc{CausalQD} are promising:

\begin{enumerate}[leftmargin=*]
    \item \textbf{Learned descriptors (AURORA-style)}: Automatically discover behavioural descriptors using autoencoders trained on the archive population, removing the need for hand-designed features.
    
    \item \textbf{Continuous relaxation}: Combine MAP-Elites with continuous DAG parameterisations (e.g., NOTEARS acyclicity constraint) for hybrid discrete-continuous search.
    
    \item \textbf{CMA-ME for DAGs}: Replace uniform random mutation with Covariance Matrix Adaptation (CMA-ME) for more directed exploration of the descriptor space.
    
    \item \textbf{Multi-objective QD}: Extend the archive to support Pareto-optimal fronts across multiple score functions (e.g., BIC vs.\ BDeu).
    
    \item \textbf{Active learning}: Use the archive to design interventional experiments that maximally reduce MEC uncertainty.
    
    \item \textbf{Large-scale parallelism}: Distribute offspring evaluation across GPU or cluster nodes for scaling to $p > 100$.
    
    \item \textbf{Non-Gaussian extensions}: Integrate CAM scores \citep{buhlmann2014cam} or nonparametric CI tests for non-linear, non-Gaussian data.
    
    \item \textbf{Time-series causal discovery}: Extend CausalQD to temporal data with dynamic Bayesian networks, using temporal descriptors (lag structure, Granger causality features).
    
    \item \textbf{Latent variables}: Handle latent confounders by extending the search space to include MAGs (maximal ancestral graphs) and PAGs (partial ancestral graphs).
    
    \item \textbf{Federated causal discovery}: Apply CausalQD in federated settings where data is distributed across sites, combining local archives via archive merging protocols.
\end{enumerate}

\end{document}
